\documentclass{article}
\usepackage[utf8]{inputenc}
\usepackage[portuguese]{babel}
\usepackage{amsmath}
\usepackage{amsfonts}
\usepackage{amssymb}
\usepackage{graphicx}
\usepackage{longtable}
\usepackage{booktabs}

\title{Documento Escaneado}
\author{}
\date{}

\begin{document}
\maketitle

\section*{3 IDENTIDADE INSTITUCIONAL DA GERÊNCIA DE TECNOLOGIA DA INFORMAÇÃO}

Identidade é tudo que torna algo único. No caso de uma equipe, podemos entender sua identidade institucional como o conjunto de suas características próprias e exclusivas, refletindo seu papel.

\subsection*{3.1 Missão da Gerência de Tecnologia da Informação}

Suprir a Fundep de soluções tecnológicas informacionais para viabilizar o alcance dos seus objetivos, fornecendo para os gestores internos e parceiros as ferramentas necessárias ao cumprimento de suas tarefas, trazendo inovação, segurança e agilidade.

\subsection*{3.2 Atividade da Gerência de Tecnologia da Informação}

Tecnologia da Informação.

\subsection*{3.3 O que não é atividade da Gerência de Tecnologia da Informação?}

Este tópico resume-se a questões que, atualmente, fazem parte do portfólio de atividades desempenhadas pela equipe, mas que apresentam natureza diferente da atividade fim, atribuindo demanda de mão de obra que retiram o foco de atividades mais específicas.

\begin{longtable}{|c|p{5cm}|p{8cm}|}
\toprule
Nº & Item & Motivo \\
\midrule
01 & Definir regras e processos para desenvolvimento de sistemas. & Gera risco institucional e retrabalho. A definição deve ser realizada pelo gestor ou solicitante. \\
02 & Atender auditorias e fiscais externos. & Gera risco institucional. É tarefa da Auditoria. \\
03 & Induzir conceitos para extração de dados. & Gera risco institucional e a responsabilidade por definir os conceitos é do solicitante. \\
04 & Manter o ambiente de sistemas de outras instituições em funcionamento (Backup, importação de dados, Pick, etc.). & Responsabilidade é da própria fundação. \\
05 & Inserir estagiário no processo de trabalho. & Responsabilidade do gestor. \\
06 & Mascarar problemas operacionais por meio de simulações de sistemas. & Risco institucional. \\
07 & Importar produtos não desenvolvidos internamente. & Impossibilidade de atendimento (inviabilidade técnica). \\
08 & Realizar manutenção e configuração em equipamentos que não pertencem à Fundep. & Risco alto e geração de alta demanda. \\
09 & Instalar equipamentos que não são inerentes à informática ou que estejam disponíveis para reserva. & Responsabilidade da Geinfra. \\
10 & Prestar consultoria para instalação de equipamentos de informática (layout). & Responsabilidade da Geinfra. \\
11 & Efetuar processos de compras. & Responsabilidade da Gecom. \\
12 & Dar suporte e monitorar a infraestrutura de terceiros (HRTN, UPA e etc). & Responsabilidade da Geinfra. \\
13 & Esta infraestrutura não pertence à Fundep. & \\
\bottomrule
\caption{O que não é atividade da Getin}
\label{tab:nao_atividade_getin}
\end{longtable}

\subsection*{3.4 O que não é atividade da Getin, mas deveria ser?}

Aqui são apresentadas atividades consideradas importantes pela equipe, mas que não fazem parte de seu portfólio atual.

\begin{longtable}{|c|p{6cm}|p{8cm}|}
\toprule
Nº & Item & Motivo \\
\midrule
01 & Dar suporte a produto padrão de mercado (Word, Excel, etc.) não desenvolvidos pela Fundep. & Disseminar informação. \\
02 & Treinar os usuários nas políticas e práticas de segurança da informação. & Instruir o usuário quanto às políticas de segurança. \\
\bottomrule
\caption{O que não é atividade da Getin, mas deveria ser?}
\label{tab:nao_atividade_getin_deveria_ser}
\end{longtable}

\subsection*{3.5 Fornecedores e Insumos}

Entende-se por fornecedores, quem abastece ou propicia os recursos necessários (matérias primas, energia, equipamentos, informações, etc.) para produzir mercadorias ou serviços. Podem ser internos ou externos.

Insumo é tudo aquilo que é transformado, modificado ou tratado na execução de um processo (matéria prima, energia, informações, uso de equipamentos, capital, horas de trabalho, etc.) necessário para produzir mercadorias ou serviços.

Assim, segue a descrição dos fornecedores e insumos relacionados.

\begin{longtable}{|c|p{4cm}|p{8cm}|}
\toprule
Nº & Fornecedores & Insumos \\
\midrule
01 & Usuários, Servidor e Getin. & Demandas (Chamados, manutenção corretiva, evolutiva, extração de dados, etc.), estabelecimento e análise de impacto das solicitações. \\
02 & Assessoria de Planejamento (SGI). & Análise de demandas. \\
03 & Diretoria & Fornecimento de demandas. \\
04 & Empresas fornecedoras de ferramenta (energia elétrica, Deltatronio, etc.) & Ferramental (Kit de ferramentas, cabos, conectores, mídias, etc.), energia. \\
05 & Governo. & Informações de servidores públicos. \\
06 & Fundep. & Layouts para arquivos e regras. \\
07 & Bancos. & Informações, críticas e identificação de erros. \\
08 & Coordenadores. & Demandas. \\
09 & Fornecedores de Hardware e Software (Dell, Printer Tec, HP, Citrix, Microsoft, Symantec, Papercut, DCS, IBM, empresa certificadora, SOS Monitores e aliados) & Hardware e software (computadores, impressoras, peças para manutenção, SLA de hardware e software, software de treinamento, sistemas operacionais, gerenciamento dos parques de impressoras e computadores, antivírus, serviços, softwares especializados, certificados digitais e manutenção de monitores). \\
10 & Prestadores de serviços em manutenção de equipamentos (nobreaks) & Manutenção preventiva e corretiva nos equipamentos. \\
11 & Sabaster Gerador. & Manutenção em equipamentos fornecedores de energia elétrica. \\
12 & Geinfra & Conferência periódica dos geradores, incluindo combustível. \\
13 & Cecom/RNP & Link de dados e voz. \\
\bottomrule
\caption{Fornecedores e Insumos}
\label{tab:fornecedores_insumos}
\end{longtable}

\subsection*{3.6 Processos}

Processo é o conjunto de atividades realizadas na geração de resultados para os parceiros e beneficiados desde o início do pedido até a entrega do produto, ou seja, processos são fluxos de trabalho que recebem um insumo (uma matéria prima, uma informação, etc) como entrada e devolve um produto ou serviço na saída, tendo início e fim bem definidos.

Os processos da Getin são:
\begin{enumerate}
    \item Atendimento aos usuários para esclarecimento de dúvidas, extração de dados ou manutenção corretiva em sistemas;
    \item Desenvolvimento de sistemas;
    \item Suporte técnico para manutenção de computadores e periféricos;
    \item Administração do Centro de Processamento de Dados Fundep;
    \item Controle de Licenciamento de Software;
    \item Controle de backup;
    \item Monitoramento de serviços de Tecnologia da Informação;
    \item Controle de rede de dados;
    \item Gerenciamento de sistemas de segurança.
\end{enumerate}

\subsection*{3.7 Pessoas}

A Getin conta com a seguinte estrutura funcional à disposição para o total desempenho das atividades:

\begin{longtable}{|c|p{6cm}|c|}
\toprule
Nº & Função & Quantidade \\
\midrule
01 & Analista de Sistemas & 01 \\
02 & Assistente Administrativo & 02 \\
03 & Estagiários & 02 \\
04 & Prestadores de Serviço & 01 \\
\bottomrule
\caption{Pessoas}
\label{tab:pessoas}
\end{longtable}

\subsection*{3.11.1 Forças}

São características positivas de destaque, no setor, que favorecem o cumprimento de sua missão (vantagem operacional).
Forças identificadas na Getin:

\begin{enumerate}
    \item Ferramentas e hardware que agilizam o processo de backup (exceto para Pick, que possui ferramenta própria);
    \item Rede paralela, velocidade e capacidade de armazenamento dos \textit{backups};
    \item No-breaks novos;
    \item Monitoramento em tempo real dos sistemas e serviços;
    \item Monitoramento do acesso à internet (Política de Recursos de Informática);
    \item Alto nível técnico da equipe;
    \item Qualidade do \textit{data center};
    \item \textit{Backup} sistemas Pick;
    \item Volume de entregas;
    \item Capacidade da equipe para prestar assessoria a processos;
    \item Flexibilidade e capacidade de adaptação da equipe;
    \item Equipamentos utilizados no desempenho das atividades cotidianas;
    \item Pró atividade da equipe;
    \item Bom conhecimento dos problemas;
    \item Ferramenta para acompanhamento e análise dos atendimentos e das demandas;
    \item Equipe enxuta;
    \item Responsabilidade e comprometimento dos membros da equipe;
    \item Ambiente físico adequado (bancada, ferramentas, etc.);
    \item Conhecimento especializado em impressoras;
    \item Boa troca de conhecimento entre a equipe;
    \item Cabeamento estruturado do prédio.
\end{enumerate}

\subsection*{3.11.2 Fraquezas}

São características negativas, no setor, que prejudicam o cumprimento de sua missão (desvantagem operacional).
Fraquezas identificadas na Getin:

\begin{enumerate}
    \item Ausência de ferramenta para \textit{backup} do Pick;
    \item Ausência de políticas de validação de \textit{backup};
    \item Local inadequado para armazenamento físico das mídias de \textit{backup};
    \item Não possui mídias distintas (fita e disco) para execução de todos os \textit{backups};
    \item Ausência de \textit{link} redundante, gerando indisponibilidade da conexão à internet;
    \item Ausência de análise periódica de segurança dos sistemas;
    \item Ausência de validação de senha (forte) em sistemas externos;
    \item Baixa qualidade da criptografia de dados na rede sem fio;
    \item Ausência de HTTPS (certificado digital);
    \item Falta de conhecimento dos demais colaboradores da Getin sobre as atividades relacionadas à infraestrutura;
    \item Falta gerenciamento das estatísticas de acesso aos sistemas e sites;
    \item Fundep não possui estrutura redundante de 100\% dos servidores e sistemas;
    \item Falta atualização do conhecimento da equipe em novas tecnologias;
    \item Necessidade de maior troca de conhecimento entre membros da Getin;
    \item Não sair do ambiente de trabalho para discutir problemas e soluções;
    \item Falta sistema anti-incêndio no \textit{data center};
    \item Baixa frequência de atualização da política de informática;
    \item Desnível técnico da equipe;
    \item Falta de processo de qualidade no desenvolvimento de \textit{software};
    \item Demora em restaurar \textit{Backup} SQL e arquivos;
    \item Atendimento de chamados de forma paliativa, sem fornecer solução definitiva;
    \item Falta de métodos e ferramentas para a geração de indicadores para Fundep;
    \item Não há ninguém pensando em inovação;
    \item Não há ninguém pensando em inovação;
    \item Equipe não realiza reuniões de alinhamentos periódicos;
    \item Baixa qualidade do registro dos chamados (retorno para os usuários e documentação do problema e solução);
    \item Programas fora do TFS e em versões diferentes de ferramenta;
    \item Espaço em disco nos servidores em produção;
    \item Volume de \textit{bugs} por sistema;
    \item Falta de rastreabilidade de atendimento de demandas e análise de impacto de negócio;
    \item Falta de local adequado para máquinas (estoque);
    \item Burocracia para doação/descarte de equipamentos;
    \item Falha na comunicação interna entre as atividades de suporte e infraestrutura;
    \item Não existe filtro adequado de ligações;
    \item Acesso irrestrito de usuários a sala da Getin;
    \item Dificuldade para atendimento dos prazos estabelecidos;
    \item Ausência de estrutura para vídeo conferência.
\end{enumerate}

\subsection*{3.11.2 Ambiente externo}

Tem por finalidade verificar as oportunidades e ameaças que estão no ambiente em que o setor está inserido, a fim de identificar as melhores maneiras de evitar e/ou usufruir dessas situações.

\subsubsection*{3.11.2.1 Oportunidades}

São forças ambientais incontroláveis pelo setor, que podem favorecer a sua ação estratégica, desde que conhecidas e aproveitadas satisfatoriamente.
Oportunidades identificadas na Getin:

\begin{enumerate}
    \item Participação em eventos, trazendo novas tecnologias;
    \item Constante redução dos custos de link de dados facilita a aquisição de \textit{link} redundante;
    \item Realizar palestras de orientação de boas práticas de uso da informática com foco em segurança;
    \item Visão da diretoria da Getin como área estratégica;
    \item Relacionamento com a UFMG (monitorar e trocar conhecimento);
    \item Orçamentos aprovados pela Diretoria para investimentos na atualização constante do parque tecnológico;
    \item Ampliação da pratica de virtualização de servidores;
    \item Possibilidade de criação do Comitê de Priorização;
    \item Computação em nuvem;
    \item Mobilidade;
    \item BI e \textit{Data Mining};
    \item Refatoração arquitetural;
    \item Treinamentos internos técnico da equipe com conhecimento da própria equipe;
    \item Estar no dia a dia do usuário;
    \item Implantação de todos os sistemas do GPF.NET, viabilizando melhorias nas interações (implantação do Financeiro/Gerenciador de Projetos);
    \item Novos \textit{softwares} e \textit{hardwares} que reduzem o volume de atendimentos;
    \item Alta disponibilidade de fornecedores;
    \item Disposição da Diretoria para montagem do ambiente físico para vídeo conferência.
\end{enumerate}

\subsubsection*{3.11.2.2 Ameaças}

São forças ambientais incontroláveis pelo setor, que criam obstáculos à sua ação estratégica, mas que poderão ou não ser evitados desde que conhecidas em tempo hábil.
Ameaças identificadas na Getin:

\begin{enumerate}
    \item Falta de orientação dos usuários quanto às boas práticas do uso da informática;
    \item Falta de monitoramento do funcionamento adequado das câmeras de segurança;
    \item Dependência externa para \textit{link} de dados (Cecom);
    \item Troca da diretoria de 04 em 04 anos;
    \item Mudança de estrutura pode gerar perda de dados;
    \item Pressão externa, gerando quebra de padrões e prioridades;
    \item Restrições orçamentárias;
    \item Volume de chamados repetitivos e desnecessários;
    \item Higienização do ambiente físico (fungos, poeira, etc.);
    \item Atendimento a hospitais (manter ambiente) e outras fundações (sistema, código específico);
    \item Desnível de conhecimento dos usuários em ferramentas que utilizam;
    \item Aquecimento do mercado de trabalho;
    \item Propostas fechadas com dependência da informática, sem consulta à Getin e outras áreas;
    \item Riscos inerentes a implantação do Sistema Financeiro;
    \item Queda de serviço (energia elétrica / internet);
    \item Baixo conhecimento técnico da GEF/Infraestrutura para montagem de equipamentos;
    \item Uso de equipamentos particulares na rede, oferecendo risco de vírus, falha de segurança da informação, etc.;
    \item Equipamentos adquiridos por taxa indireta não podem seguir o padrão Fundep;
    \item Falta de planejamento nas mudanças físicas;
    \item As estruturas das janelas oferecem risco constante de dano aos equipamentos.
\end{enumerate}

\subsection*{3.12 Objetivos e ações}

Os objetivos sugeridos pela equipe, para que fossem considerados válidos, foram analisados para verificar se são coerentes, viáveis, porém desafiantes; aprazados, mensuráveis, claros, explícitos e concisos; conhecidos e acreditados por toda a equipe, em número reduzido para evitar a dispersão. A seguir a lista de objetivos finais, definidos para a equipe Getin.

\begin{longtable}{|c|p{5cm}|c|p{8cm}|}
\toprule
Nº & Objetivo & Nº & Ação \\
\midrule
01 & Implantar, atendendo aos prazos e com qualidade, o sistema financeiro, até abril de 2012. & 01 & Documentar decisões e fatos importantes. \\
 & & 02 & Reforçar as ações da TI. \\
 & & 03 & Comparar dados de atendimento. \\
02 & Realizar a matriz de acordo com o usuário na Getin, até abril de 2012. & 01 & Divulgar restrição de acesso para os usuários. \\
 & & 02 & Negociar com a Diretoria local seguro para armazenamento dos backups (cofre). \\
 & & 03 & Criar e documentar procedimento de conferência e restore de backups. \\
 & & 04 & Definir responsáveis por ambientes para testes de backups. \\
03 & Realizar backup de acordo com critério definido, até março de 2012. & 01 & Criação de ambientes de testes. \\
 & & 02 & Utilizar mídias diferentes para backup. \\
 & & 03 & Atualizar softwares Test, utilizando duplicação Atualizada. \\
04 & Disponibilizar redundância de link até abril de 2012. & 01 & Acompanhar e facilitar a implantação (configurar roteadores, configurando firewall, DNS, testes e implantação). \\
05 & Formatar projeto de redundância de servidores e serviços, até julho de 2012. & 01 & Formatar projeto de redundância interna (SQL Server, Aplicação, Pick, AD, Firewall) e local. \\
06 & Melhorar o tempo e a qualidade de atendimento dos chamados, até julho de 2012. & 01 & Realizar análise de soluções paliativas e existentes dos sistemas. \\
 & & 02 & Apresentar resultado da análise e plano de ação. \\
 & & 03 & Acompanhar a quantidade de chamados corretivos de sistemas da nova plataforma em 20\%. \\
 & & 04 & Comparar dados de atendimento (dez-fev) com (mar-jul), até abril de 2012. \\
 & & 05 & Efetuar substituição de impressoras previstas em orçamento. \\
07 & Melhorar qualidade das entregas, até julho de 2012. & 01 & Realizar manutenção preventiva das impressoras conforme cronograma (Anexo II). \\
 & & 02 & Reduzir em 50\% os chamados referentes hardware de impressoras. \\
 & & 03 & Atender chamados de INFRA em até 02 dias após com taxa de retorno máximo de 2\%. \\
 & & 04 & Atender chamados de SUPORTE em até 02 dias após com taxa de retorno máximo de 2\%. \\
 & & 05 & Atender >95\% dos chamados de extração de dados e manutenção corretiva e atendimento em até 2 dias úteis. \\
 & & 06 & Definir responsáveis dos módulos/fluxo de aprovação. \\
08 & Melhorar a segurança da informação de acordo com norma 27001 03 ABNT, até julho de 2012. & 01 & Criar ambiente de desenvolvimento, homologação e pré-produção. \\
 & & 02 & Criar área de qualidade (processos e pessoal), considerando investimentos em testes unitários, padronização de comentários de código. \\
 & & 03 & Criar área de qualidade (processos e pessoas), considerando investimentos em testes unitários, padronização de comentários de código. \\
 & & 04 & Apresentar plano para prevenção e combate à incêndio para Diretoria. \\
 & & 05 & Implementar política de segurança para sistemas externos (para a equipe de desenvolvimento), até março de 2012. \\
 & & 06 & Implementar política de senha forte, alinhando junto à equipe de desenvolvimento, formando a melhor proposta, oferecendo o suporte necessário. \\
09 & Promover o nivelamento técnico da equipe, até julho de 2012. & 01 & Documentar o conhecimento específico nos pares e documentar (fora/sistema). \\
 & & 02 & Workshops internos. \\
 & & 03 & Promover treinamentos em Exchange, Windows Server, Tivoli e Symantec. \\
10 & Atingir 99,5\% de disponibilidade da infraestrutura, até julho de 2012. & 01 & Implementar política de segurança para sistemas externos (para a equipe de desenvolvimento), até março de 2012. \\
11 & Realizar auditoria semestral de licenciamento de softwares, até julho de 2012. & 01 & Conforme cronograma específico (Anexo I). \\
12 & Atualizar 100\% dos sistemas operacionais das estações de trabalho de usuários para Windows 7, até julho de 2012. & 01 & Realizar levantamento para definição de empresas/locais/entidades para recebimento e descarte em conjunto com GEF/Contabilidade. \\
13 & Manter no máximo 10 equipamentos aguardando doação/descarte, sendo que o processo para tal não ultrapasse 2 semanas, até julho de 2012. & 01 & Negociar e alinhar com a GEF/Infraestrutura o local definitivo para armazenamento e depósito de materiais de informática. \\
 & & 02 & Implementar o controle de evolutivas para 100\% das entregas. Acrescentar a evolutiva no controle em no máximo 5 dias úteis. \\
14 & Implementar cubo financeiro, folha e gerência, até julho de 2012. & 01 & Criar documento oficial para envio formalizado. \\
15 & Relacionar responsabilidades para clientes externos, até maio de 2012. & 01 & Implementar política de segurança para sistemas externos (para a equipe de desenvolvimento), até março de 2012. \\
16 & Implementar política de segurança para sistemas externos (para a equipe de desenvolvimento), até março de 2012. & 01 & Implementar política de segurança para sistemas externos (para a equipe de desenvolvimento), até março de 2012. \\
\bottomrule
\caption{Objetivos e ações}
\label{tab:objetivos_acoes}
\end{longtable}

\end{document}
